\chapter{릴레이와 접점 논리}\label{ch:relay}

1--4부는 프로그램과 디지털 회로의 코드였다. 5부는 전시의 두 번째 기계, 릴레이식 봄베의 회로를 다룬다.
릴레이식 봄베에는 FPGA도, 트랜지스터로 된 논리 게이트도 거의 없다. 계산을 하는 것은 전자석이 끌어당기는
금속 접점이다. 이 장은 릴레이 한 개가 어떻게 생겼는지에서 출발해, 접점을 엮어 논리를 만드는 법, 그리고
그 방법으로 지은 네 가지 회로(위치 표시의 접점 나무, 이진 시프터, 오르기만 하는 계수기)를 차례로 본다.
전기를 몰라도 읽을 수 있게 썼다. 필요한 것은 ``전류는 끊기지 않은 길로만 흐른다''는 사실 하나다(전압·전류·다이오드가 낯설면
\ref{ch:39-electric-basics}장을 먼저 본다).

\section{릴레이 한 개: 코일과 접점}\label{sec:relay-one}

\Why \term{릴레이}{relay}는 전기로 움직이는 스위치다. 철심에 감은 전선(\term{코일}{coil})에 전류가 흐르면
철심이 자석이 되어 얇은 철판(\term{접극자}{armature})을 끌어당긴다. 접극자에는 금속 접점이 붙어 있어서,
끌려가면 접점이 한쪽 단자에서 떨어져 다른 쪽 단자에 닿는다. 전류를 끊으면 스프링이 접극자를 제자리로
돌려보낸다(그림~\ref{fig:r-anatomy}). 코일 쪽 회로와 접점 쪽 회로는 전기적으로 이어져 있지 않다. 작은 코일
전류 하나로 전혀 다른 회로의 길을 열고 닫을 수 있다는 것이 릴레이의 쓸모다.

접점에는 이름이 셋 있다. 움직이는 쪽을 \term{공통}{COM}이라 하고, 코일이 꺼져 있을 때 공통과 닿아 있는
단자를 \term{상시 닫힘}{NC, normally closed}, 코일이 켜져야 닿는 단자를 \term{상시 열림}{NO, normally open}이라
한다. 우리말 문헌에서는 NO를 a접점, NC를 b접점이라고도 부른다. 공통 하나에 NC와 NO가 모두 있는 접점을
\term{전환 접점}{changeover, c접점}이라 한다. 코일 하나가 전환 접점을 몇 벌 움직이는지를 \term{극}{pole}의
수로 센다. 이 프로젝트의 릴레이는 거의 전부 극이 둘인 전환 접점 릴레이, 곧 \term{DPDT}{double pole, double
throw}다.

\begin{figure}[ht]
\centering
\begin{tikzpicture}[x=1cm,y=1cm]
  \foreach \dx/\on/\ttl in {0/0/{코일 꺼짐},7/1/{코일 켜짐}} {
    \begin{scope}[xshift=\dx cm]
      \node[font=\small\bfseries] at (2.2,3.3) {\ttl};
      % core and coil
      \fill[cgray!40] (0,0) rectangle (0.6,2.0);
      \foreach \y in {0.25,0.5,...,1.75} \draw[corange,line width=1.2pt] (-0.1,\y) -- (0.7,\y);
      \node[note,anchor=north] at (0.3,-0.1) {코일};
      \ifnum\on=1
        \foreach \y in {0.5,1.0,1.5} \draw[->,cblue] (0.9,\y) -- (1.25,\y);
      \fi
      % armature (pivot at (0.3,2.3))
      \fill (0.3,2.35) circle (0.06);
      \ifnum\on=0
        \draw[line width=2pt,cgray] (0.3,2.35) -- (3.2,2.75);
        \draw[line width=1pt] (3.2,2.75) -- (3.2,2.35);
        \fill[cblue] (3.2,2.35) circle (0.07);
      \else
        \draw[line width=2pt,cgray] (0.3,2.35) -- (3.2,2.05);
        \draw[line width=1pt] (3.2,2.05) -- (3.2,1.45);
        \fill[cblue] (3.2,1.45) circle (0.07);
      \fi
      % NC (top) and NO (bottom) terminals
      \draw[line width=1pt] (3.0,2.42) -- (4.2,2.42) node[right,font=\footnotesize] {NC};
      \draw[line width=1pt] (3.0,1.38) -- (4.2,1.38) node[right,font=\footnotesize] {NO};
      \draw[line width=1pt] (0.3,2.35) -- (0.3,2.9) -- (-0.4,2.9) node[left,font=\footnotesize] {COM};
      \draw[cgreen] (3.9,2.6) -- (3.8,2.68) -- (4.0,2.76) -- (3.8,2.84) -- (4.0,2.92) -- (3.8,3.0) -- (3.9,3.08);
      \node[note,anchor=west] at (4.0,3.0) {스프링};
      \ifnum\on=0
        \node[note,text=cblue,anchor=west] at (0.9,0.7) {COM--NC가\\닿아 있다};
      \else
        \node[note,text=cblue,anchor=west] at (1.35,0.7) {COM--NO가\\닿아 있다};
      \fi
    \end{scope}
  }
\end{tikzpicture}
\caption{릴레이의 동작. 코일에 전류가 흐르면 철심이 접극자를 끌어당겨 공통(COM)이 NC에서 떨어지고 NO에 닿는다.
코일 전류를 끊으면 스프링이 되돌린다.}
\label{fig:r-anatomy}
\end{figure}

릴레이는 느리다. 코일에 전압을 건 뒤 접점이 실제로 넘어가 흔들림(\term{튐}{bounce})까지 가라앉는 데 이
프로젝트의 소형 신호 릴레이로 2.5--5~ms가 걸리고, 놓일 때도 몇 ms가 걸린다. FPGA의 플립플롭이 수 나노초에
바뀌는 것과 비교하면 백만 배쯤 느리다. 이 느림과, 같은 품번이라도 개체마다 시간이 다르다는 점이
\ref{ch:relay-timing}장의 주제다.

\Ours 릴레이식 봄베의 모든 보드는 Omron G6S-2급의 24~V DPDT 신호 릴레이 한 종류를 쓴다. 회로도 생성기는
핀 번호와 기능의 대응을 한 줄로 적어 두고 모든 보드가 이 줄을 공유한다.

\begin{lstlisting}[language=Python,caption={\code{gen\_shifter.py}: G6S-2의 핀 기능. 코일은 1(+)과 12($-$), 극 1은 4/3/5, 극 2는 9/10/8.},label={lst:r-pins}]
# G6S-2 pin functions (verified against the library symbol graphics: the armature of pole 1
# rests on pin 3, of pole 2 on pin 10, when the coil is off).
RELAY_PINS = dict(coil_p="1", coil_n="12", c1="4", nc1="3", no1="5", c2="9", nc2="10", no2="8")
\end{lstlisting}

검증기도 같은 대응을 따로 적는다(\code{POLES = [("4", "3", "5"), ("9", "10", "8")]}). 생성기와 검증기가 같은
사실을 각자 적어 두는 이유는 \ref{ch:boardverify}장에서 본다.

\section{접점으로 논리를 만든다}\label{sec:relay-logic}

\Why \ref{ch:hdl}장의 게이트는 0과 1을 전압으로 나타냈다. 릴레이 회로에서도 같다. 어떤 선에 +24~V가 닿아
있으면 1, 아니면 0이다. 다만 계산하는 방법이 다르다. 게이트는 입력 전압을 보고 출력 전압을 새로 만들지만,
접점은 길을 열고 닫을 뿐이다. 그래서 접점 회로의 논리는 ``전원에서 출발해 닫힌 접점만 밟아 이 점에 올 수
있는가''라는 길 찾기 문제가 된다. 이 관점은 \ref{ch:boardverify}장의 시뮬레이터가 그대로 쓴다.

\begin{itemize}
\item 두 접점을 \emph{직렬}로 이으면 둘 다 닫혀야 길이 생긴다. AND다.
\item 두 접점을 \emph{병렬}로 이으면 하나만 닫혀도 길이 생긴다. OR다.
\item NC 접점은 코일이 꺼져 있을 때 닫힌다. NOT이다.
\end{itemize}

\General 계단 조명의 3로 스위치, 자동차의 경적 릴레이, 엘리베이터의 층 선택 회로는 모두 이 세 규칙으로
만든 것이다. 그림~\ref{fig:r-logic}는 램프 하나를 세 가지 방식으로 켜는 회로다.

\begin{figure}[ht]
\centering
\begin{circuitikz}[x=1cm,y=1cm]
  \ctikzset{bipoles/length=9mm}
  \foreach \dx/\ttl in {0/AND (직렬),5/OR (병렬),10/NOT (NC 접점)} \node[font=\small\bfseries] at (\dx+1.7,2.4) {\ttl};
  % AND
  \draw (0,1.6) node[left,font=\footnotesize] {$+24$ V} to[nos,l={\scriptsize $a$}] (1.4,1.6) to[nos,l={\scriptsize $b$}] (2.8,1.6) -- (3.3,1.6) to[lamp] (3.3,0) -- (0,0) node[left,font=\footnotesize] {0 V};
  % OR
  \begin{scope}[xshift=5cm]
  \draw (0,1.6) node[left,font=\footnotesize] {$+24$ V} -- (0.4,1.6) coordinate (s);
  \draw (s) to[nos,l={\scriptsize $a$}] (2.2,1.6) coordinate (e);
  \draw (s) -- (0.4,0.9) to[nos,l_={\scriptsize $b$}] (2.2,0.9) -- (e);
  \draw (e) -- (3.3,1.6) to[lamp] (3.3,0) -- (0,0) node[left,font=\footnotesize] {0 V};
  \end{scope}
  % NOT
  \begin{scope}[xshift=10cm]
  \draw (0,1.6) node[left,font=\footnotesize] {$+24$ V} to[ncs,l={\scriptsize $a$}] (2.2,1.6) -- (3.3,1.6) to[lamp] (3.3,0) -- (0,0) node[left,font=\footnotesize] {0 V};
  \end{scope}
\end{circuitikz}
\caption{접점 논리. 접점 $a$, $b$는 각각 다른 릴레이의 접점이다(소문자는 접점, 그 릴레이의 코일은 대문자로
적는다). 열린 스위치 기호가 NO, 닫힌 기호가 NC다.}
\label{fig:r-logic}
\end{figure}

하나 더 있다. \term{다이오드}{diode}는 전류를 한 방향으로만 흘리는 부품이다. 여러 신호를 다이오드 하나씩을
거쳐 한 점에 모으면, 어느 하나라도 +24~V이면 그 점이 +24~V가 되고 거꾸로 그 점의 전압이 다른 입력으로
새어 나가지 않는다. 이것을 \term{다이오드 OR}{diode OR}라고 한다. 접점 하나를 쓰지 않고 OR를 만들 수 있어
이 프로젝트의 계수기는 다이오드를 자리마다 64개나 쓴다.

\Ours 접점 논리가 게이트 논리와 가장 다른 점은 \emph{양방향}이라는 것이다. 게이트는 입력에서 출력으로만
신호를 보내지만, 닫힌 접점은 그냥 전선이라 양쪽으로 전류가 흐른다. 그래서 의도하지 않은 우회로(\term{샛길}{sneak
path})가 생기기 쉽다. 이 프로젝트의 시프터 보드는 ``한 단의 릴레이가 모두 함께 움직인다''는 성질로
샛길을 막고(\ref{sec:relay-shifter}절), 테스트 레지스터는 선마다 다이오드를 넣어 막는다(\ref{sec:t-register}절).
검증기가 ``기대한 연결 \emph{말고는} 아무것도 이어지지 않았는가''를 꼭 확인하는 이유도 이것이다.

\section{자기 유지: 릴레이로 만든 기억}\label{sec:relay-latch}

\Why 릴레이의 접점이 자기 코일의 전원 길에 들어가면 기억이 생긴다. 한 번 켜진 릴레이가 자기 NO 접점으로
자기 코일에 전류를 계속 대 주기 때문이다. 처음 켠 신호가 사라져도 릴레이는 켜진 채 남는다. 끄려면 그
자기 유지 길 어딘가를 따로 끊어야 한다. 이것이 \term{자기 유지}{self-hold} 회로이고, \ref{ch:sync}장의
플립플롭이 하던 ``1비트 기억''을 릴레이 하나와 접점 하나로 만든 것이다.

\General 공장 기계의 시작/정지 버튼이 전형적인 예다(그림~\ref{fig:r-latch}). 시작 버튼(NO)을 누르면
릴레이 $K$가 켜지고, $K$의 접점 $k$가 시작 버튼과 병렬로 닫혀 버튼을 놓아도 $K$가 남는다. 정지 버튼(NC)은
이 길 전체와 직렬이어서, 누르면 길이 끊겨 $K$가 놓이고 $k$도 열린다. 정전 뒤에 기계가 저절로 다시 돌지
않는다는 안전상의 이점도 있다.

\begin{figure}[ht]
\centering
\begin{circuitikz}[x=1cm,y=1cm]
  \ctikzset{bipoles/length=9mm}
  \draw (0,1.6) node[left,font=\footnotesize] {$+24$ V} to[ncs,l={\scriptsize 정지}] (1.6,1.6) coordinate (a)
        to[nos,l={\scriptsize 시작}] (3.4,1.6) coordinate (b) -- (4.2,1.6);
  \draw (a) -- (1.6,0.7) to[nos,l_={\scriptsize $k$ (자기 유지)}] (3.4,0.7) -- (b);
  \node[rcoil,anchor=west] (K) at (4.2,1.6) {$K$};
  \draw (K.east) -- ++(0.4,0) node[right,font=\footnotesize] {0 V};
  \node[circ] at (a) {}; \node[circ] at (b) {};
\end{circuitikz}
\hspace{6mm}
\begin{tikztimingtable}[timing/slope=0.1,timing/coldist=4pt,xscale=1.5,yscale=1.2,semithick]
  시작 & 1L .6H 4.4L \\
  정지 & 4L .6H 1.4L \\
  $K$   & 1.2L 2.9H 1.9L \\
\end{tikztimingtable}
\caption{시작/정지 회로. 시작 버튼을 잠깐 눌러도 $K$는 자기 접점 $k$로 켜진 채 남고, 정지 버튼이 길을 끊어야
놓인다. 오른쪽의 $K$가 버튼보다 조금 늦게 바뀌는 것은 릴레이의 동작 시간이다.}
\label{fig:r-latch}
\end{figure}

\Ours 자기 유지는 이 프로젝트의 릴레이 회로 곳곳에 있다.
\begin{itemize}
\item 테스트 레지스터의 감지 릴레이 26개는 시험 전압이 닿은 짧은 순간에 붙고, 자기 유지로 그 결과를 붙든
뒤 판정을 받는다(\ref{sec:t-register}절). 유지 길을 공통의 \term{유지 버스}{hold bus}에서 먹이므로 그 버스
하나를 끊으면 26개가 한꺼번에 지워진다.
\item 계수기의 슬레이브 릴레이는 자기 유지로 값을 기억한다(\ref{sec:relay-counter}절).
\item 정지 릴레이는 정지를 붙들어 기계를 세워 두고, 해제 버튼이 그 유지를 끊는다.
\end{itemize}

\section{먼저 끊고 나중에 잇기}\label{sec:relay-bbm}

\Why 전환 접점의 공통은 NC에서 떨어져 NO로 \emph{날아간다}. 그 사이 잠깐(이 프로젝트의 시뮬레이터는
0.2--0.6~ms로 잡는다) 공통은 어디에도 닿아 있지 않다. 이 성질을 \term{먼저 끊고 나중에 잇기}{break-before-make}라
한다. 덕분에 NC 쪽 회로와 NO 쪽 회로는 \emph{결코 동시에} 이어지지 않는다. 한 접점이 두 길 가운데 정확히
하나를 고르는 것이 기계적으로 보장되는 것이다.

릴레이 두 개로 같은 일을 하면 이 보장이 사라진다. 하나는 켜지고 하나는 꺼지는 순간에 두 릴레이의 시간이
다르면, 잠깐 둘 다 닫혀 있거나 둘 다 열려 있게 된다(그림~\ref{fig:r-bbm}).

\begin{figure}[ht]
\centering
\begin{tikztimingtable}[timing/slope=0.1,timing/coldist=4pt,xscale=1.6,yscale=1.2,semithick]
  코일 & 1L 4H 1L \\
  NC 쪽 & 1.6H 4L .4H \\
  NO 쪽 & 1.8L 3.4H .8L \\
  \textcolor{cred}{릴레이 둘: A} & 1.4H 4L .6H \\
  \textcolor{cred}{릴레이 둘: B} & 1.2L 4.6H .2L \\
\extracode
  \begin{pgfonlayer}{background}
    \fill[cgreenbg] (1.6,-5.0) rectangle (1.8,1.0);
    \fill[credbg] (1.2,-9.6) rectangle (1.4,-5.4);
    \fill[credbg] (5.4,-9.6) rectangle (5.8,-5.4);
  \end{pgfonlayer}
  \node[note,text=cgreen,anchor=west] at (6.2,-2.2) {전환 접점 하나: 끊긴 틈(초록)은\\있어도 겹침은 없다};
  \node[note,text=cred,anchor=west] at (6.2,-7.6) {따로 움직이는 두 릴레이:\\시간이 어긋나면 둘 다 닫힌다(빨강)};
\end{tikztimingtable}
\caption{먼저 끊고 나중에 잇기. 위는 전환 접점 하나의 NC와 NO, 아래는 같은 일을 릴레이 두 개(A는 NC처럼,
B는 NO처럼 쓴 것)에 맡긴 경우다.}
\label{fig:r-bbm}
\end{figure}

\General 철도의 선로 전환기가 같은 원리다. 레일이 한쪽에서 떨어진 뒤에야 다른 쪽에 붙으므로 열차가 두 선로에
동시에 이어지는 일이 없다. 전원 회로에서는 정반대 성질(먼저 잇고 나중에 끊기)이 필요할 때도 있다. 그래서
부품 자료에는 어느 쪽인지가 꼭 적혀 있다.

\Ours 릴레이식 계수기는 클럭을 \code{CK}와 \code{NCK} 두 선으로 받는다. 둘은 반드시 정반대여야 하고, 둘이
동시에 켜지면 계수기 값이 틀어진다. 그래서 순서기는 이 두 선을 \emph{릴레이 하나의 한 극}에서 낸다. 공통에
+24~V, NO가 \code{CK}, NC가 \code{NCK}다. 계수기 보드의 설명서는 ``개별 릴레이 둘로 만들어 둘 다 켜지면 값이
틀어진다(시뮬레이션에서 확인)''고 적고 있고, 순서기의 구조 검사는 \code{CK}와 \code{NCK}가 버스마다 한
릴레이의 한 극에서만 나오는지를 확인한다. 끊긴 틈 동안에는 둘 다 꺼지는데, 이 틈이 릴레이가 놓이는 시간보다
훨씬 짧아 계수기의 릴레이들은 그동안 상태를 유지한다.

\section{역기전력 다이오드: 코일을 끌 때}\label{sec:relay-flyback}

\Why 코일은 전류의 변화를 싫어한다. 코일에 흐르던 전류를 갑자기 끊으면 코일은 그 전류를 계속 흘리려고
매우 큰 전압을 만든다(\term{역기전력}{back EMF}). 전압의 크기는 전류가 줄어드는 빠르기에 비례하므로,
스위치가 순식간에 열리면 수백 V가 나올 수도 있다. 이 전압은 코일을 끈 스위치(접점이나 트랜지스터)를
상하게 하고 주변 회로에 잡음을 뿌린다.

해법은 코일 양 끝에 다이오드를 거꾸로 거는 것이다(그림~\ref{fig:r-flyback}). 평소에는 다이오드가 막혀 있어
아무 일도 하지 않는다. 스위치가 열리는 순간 코일의 전류는 다이오드를 지나 코일로 되돌아가는 작은 고리를
돌며 천천히 사라진다. 코일 양 끝의 전압은 다이오드 하나의 전압(약 0.7~V)으로 묶인다. 이 다이오드를
\term{역기전력 다이오드}{flyback diode}라 한다.

\begin{figure}[ht]
\centering
\begin{circuitikz}[x=1cm,y=1cm]
  \ctikzset{bipoles/length=9mm}
  \foreach \dx/\ttl in {0/켜져 있을 때,6.5/스위치를 연 직후} {
  \begin{scope}[xshift=\dx cm]
    \node[font=\small\bfseries] at (1.9,3.3) {\ttl};
    \draw (0,2.5) node[left,font=\footnotesize] {$+24$ V} -- (0.3,2.5);
    \ifdim\dx pt=0pt
      \draw (0.3,2.5) to[cspst] (1.6,2.5);
    \else
      \draw (0.3,2.5) to[ospst] (1.6,2.5);
    \fi
    \draw (1.6,2.5) -- (2.6,2.5) to[L,l_={코일}] (2.6,0) -- (0,0) node[left,font=\footnotesize] {0 V};
    \draw (2.6,2.5) -- (3.8,2.5);
    \draw (3.8,0) to[D*,l_={\scriptsize 다이오드}] (3.8,2.5);
    \draw (3.8,0) -- (2.6,0);
    \node[circ] at (2.6,2.5) {}; \node[circ] at (2.6,0) {};
  \end{scope}}
  \draw[->,cblue,line width=1pt] (1.2,2.8) -- (2.2,2.8);
  \draw[->,cblue,line width=1pt] (2.35,2.1) -- (2.35,0.6);
  \draw[->,corange,line width=1pt] (8.85,2.1) -- (8.85,0.5);
  \draw[->,corange,line width=1pt] (9.4,-0.25) -- (10.0,-0.25);
  \draw[->,corange,line width=1pt] (10.55,0.5) -- (10.55,2.0);
  \draw[->,corange,line width=1pt] (10.0,2.75) -- (9.4,2.75);
  \node[note,text=corange,anchor=west] at (11.8,1.3) {코일의 전류가\\다이오드로 돌며\\사라진다};
\end{circuitikz}
\caption{역기전력 다이오드. 다이오드는 평소에 막혀 있다(삼각형이 가리키는 쪽이 $+$). 스위치가 열리면 코일의
전류가 다이오드를 지나는 고리로 돌아 천천히 줄어든다.}
\label{fig:r-flyback}
\end{figure}

대가가 있다. 전류가 천천히 사라지므로 릴레이가 \emph{늦게 놓인다}. 이 프로젝트는 다이오드만 단 코일의
복귀 시간을 데이터시트 값의 2--3배인 3--10~ms로 추정한다. 빨리 놓여야 하는 릴레이에는 다이오드에 제너
다이오드(정해진 전압을 넘으면 거꾸로도 흐르는 다이오드)를 직렬로 더한다. 코일 전압이 0.7~V가 아니라
약 31~V에서 묶이므로 전류가 훨씬 빨리 줄어 복귀가 1.5--4~ms로 짧아진다. 에너지를 빨리 버리는 대신 순간
전압이 커지는 맞바꿈이다.

\General 트랜지스터로 릴레이나 모터를 켜고 끄는 회로라면 어디에나 이 다이오드가 있다. 아두이노로 릴레이를
움직이는 입문 예제에서 다이오드를 빠뜨리면 트랜지스터가 몇 번 만에 망가진다.

\Ours 이 프로젝트의 모든 릴레이 코일에는 역기전력 다이오드(1N4148W)가 하나씩 있다. 시프터 보드 한 장에
릴레이 130개와 다이오드 130개가 있는 이유다. 생성기는 릴레이를 놓을 때마다 다이오드를 같이 놓는다.

\begin{lstlisting}[language=Python,caption={\code{gen\_shifter.py}: 릴레이마다 코일 양 끝에 다이오드(음극이 $+$ 쪽).},label={lst:r-flyback}]
        sh.pin_net(S_RELAY, x, y, P["coil_p"], "CTRL")
        sh.pin_net(S_RELAY, x, y, P["coil_n"], "GND")
        # flyback diode across the coil (cathode to the + side)
        dref = f"D{st}{i + 1:02d}"
        dx = g(x - 24.0)
        sh.symbol(S_DIODE, dx, y, dref, "1N4148W", FP_DIODE, "", "Coil flyback diode", path, show_value=False)
        sh.pin_net(S_DIODE, dx, y, "1", "CTRL")   # K
        sh.pin_net(S_DIODE, dx, y, "2", "GND")    # A
\end{lstlisting}

테스트 레지스터의 타이밍 릴레이 셋(\code{STROBE}, \code{JUDGE}, \code{CLEAR})은 빨리 놓여야 하므로 다이오드 +
제너 30~V를 쓴다. 그러면 새 문제가 생긴다. 신호를 끄는 순간 코일의 $+$ 쪽이 약 $-31$~V까지 내려가는데, 그
$+$ 쪽이 곧 커넥터의 핀이면 그 음전압이 케이블을 타고 신호를 내 준 쪽(기어식의 광센서 출력 등)으로
나간다. 그래서 결정 0003은 ``받는 쪽이 역기전력을 자기 안에서 처리한다''는 규칙을 두고, 핀과 코일 사이에
다이오드를 하나 더 넣었다. 끄는 순간 이 다이오드가 막혀 코일의 전류는 보드 안의 클램프로만 돈다. 코일
전압은 다이오드 강하 0.7~V만큼 줄어든다.


\section{정리}

\begin{itemize}
\item 릴레이는 코일 전류로 접점을 옮기는 스위치다. 접점은 COM, NC(b접점), NO(a접점)이고, 이 프로젝트는 극이
둘인 전환 접점 릴레이(DPDT) 한 종류를 쓴다.
\item 접점 논리는 길 찾기다. 직렬은 AND, 병렬은 OR, NC는 NOT, 다이오드 여럿은 OR. 접점은 양방향이라 샛길에
주의한다.
\item 자기 유지 접점이 1비트 기억을 만든다. 끄려면 유지 길을 따로 끊는다.
\item 전환 접점은 먼저 끊고 나중에 잇는다. 정확히 하나를 골라야 하는 두 신호(\code{CK}/\code{NCK})는 한 접점에서
낸다.
\item 코일마다 역기전력 다이오드를 단다. 늦게 놓이는 대가가 있어, 빨리 놓여야 하면 제너를 더한다.
\end{itemize}
